\clearpage
\begingroup
\fancyhead[L]{}
\renewcommand{\headrulewidth}{0pt}
\SourceLeaf{646}
\needspace{6\baselineskip}
\pdfbookmark[1]{Bemerkungen zur Theorie der Bewegung des Mondes}{moon-bemerkungen}
\begin{center}\large BEMERKUNGEN ZUR THEORIE DER BEWEGUNG DES MONDES.\end{center}

\par\noindent\hspace*{\parindent}Die vorstehend abgedruckte »Theorie der Bewegung des Mondes«, die auf einem Konvolut von
Blättern steht, ist unvollendet geblieben. Außer dem abgedruckten Manuskript befinden sich im Nachlass noch zwei
frühere, weniger ausgeführte Entwürfe. In \textsc{Gauss}’ wissenschaftlichem Tagebuch findet sich unter 1801
August die Notiz »\textit{Theoriam motus Lunae aggressi sumus}«. Die Entstehung des Manuskripts wird man
tatsächlich in diese Zeit setzen müssen. Neben der Handschrift und der Darstellungsweise spricht auch der
Umstand für diese frühe Entstehungsperiode, dass die Anomalien vom Apogäum statt vom Perigäum gezählt
sind. Vgl. dazu die Ceresstörungen.

\par\noindent\hspace*{\parindent}Das Interesse, das diese \textsc{Gauss}schen Untersuchungen gegenüber den Methoden anderer, auch späterer
Autoren, bieten, konzentriert sich auf die ihm eigene Behandlung der Differentialgleichungen der Mondbewegung.
Die Durchführung der Integration und die Ermittlung der Störungsglieder ist unvollständig
geblieben. Nur die Breitenstörungen sind vollständig, und offenbar hat \textsc{Gauss} diese Papiere später nie
wieder zur Hand genommen. Auch hat er sich sonst nirgends darüber geäußert, auch nicht in dem 1802
beginnenden Briefwechsel mit \textsc{Olbers}. Im Jahr 1802 erschien auch bereits der dritte Band von \textsc{Laplaces}
Mécanique céleste mit der Mondtheorie, wodurch die von \textsc{Gauss} begonnenen Entwicklungen überholt
wurden. Auch dieser Umstand sowie der direkte Anschluss an \textsc{Tobias Mayer} sprechen für die erwähnte
frühe Abfassungszeit.

\par\noindent\hspace*{\parindent}Die Form, in der sich schließlich bei \textsc{Gauss} die Störungen darstellen, ist dieselbe wie die \textsc{Plana}sche
(1832), da beide die Divisoren nach Potenzen der Größe \(n\) (Verhältnis der Umlaufzeiten von Mond und
Sonne) entwickeln.

\par\noindent\hspace*{\parindent}Einige kleine Zusätze und Verbesserungen sind gemacht worden, um die Resultate wenigstens in den
von \textsc{Gauss} zunächst angenommenen Grenzen vollständig und richtig zu geben. So sind bei der Integration
der zweiten Gleichung des dritten Abschnitts, S.~622, die Glieder in \(\varepsilon\sin(2E\pm a)\) hinzugefügt und in
den folgenden Gleichungen beibehalten worden, obwohl sie, wie \textsc{Gauss} bemerkt, dort fortgelassen werden
können, weil sie auch nach den Integrationen nur von der dritten Ordnung sind. Indessen ist ihre Berücksichtigung
für die zweite Annäherung erforderlich, um die Koeffizienten der analogen Glieder im Ausdruck
für \(\frac{\mathrm du^2}{\mathrm dv^2}\pi\), S.~634, richtig zu erhalten.

\par\noindent\hspace*{\parindent}Im vierten Abschnitt ist die Differentialgleichung für \(\theta\) bis zu den Gliedern 5.~Ordnung einschließlich
genau. Nach der Integration fehlen also naturgemäß die Glieder, welche durch den Divisor \(n\) in die
5.~Ordnung aufrücken, sowie das Glied in \(\sin(2E-\lambda+2a)\), das den Divisor \(nn\) erhält und dessen
Fehlen \textsc{Gauss} durch einen \(*\) andeutet. Es ergibt sich nach \textsc{Plana} zu \(-\frac{9}{32}n\varepsilon\varepsilon\sin(2E-\lambda+2a)\).
Bei der Transformation auf die \textsc{Mayer}sche Form, S.~631---633, sind die Verzeichnisse der hinzutretenden Glieder,
welche übrigens nicht vollständig sind, beim Abdruck fortgelassen worden. Auch ist die Gegenüberstellung
der \textsc{Gauss}schen und \textsc{Mayer}schen Resultate keiner Kontrollrechnung unterworfen worden.

\VIIInlineSourceHead{\SourceLeaf{647}}{}{BEMERKUNGEN.\quad THEORIE DER BEWEGUNG DES MONDES.}{641}
\par\noindent\hspace*{\parindent}Im fünften Abschnitt hat \textsc{Gauss} die Differentialgleichung für \(p\) nur bis zu den Gliedern vierter Ordnung
einschließlich aufgestellt. Bei der Entwicklung des Ausdrucks für \(\sin 2(v-V)\) sind indessen einige
Glieder dritter Ordnung mitgenommen, weil sie bei der Integration kleine Divisoren erhalten. Die beiden
Glieder in Klammern sind hier lediglich deswegen hinzugesetzt, weil \textsc{Gauss} bei Aufstellung des darauf
folgenden Ausdrucks sie berücksichtigt hat. Der schließliche Ausdruck von \(p\) enthält selbstverständlich
die Glieder unterhalb der 5.~Ordnung nicht, die bei der Integration aus Gliedern 5.~Ordnung entstehen.
Einzelne dieser Glieder hatte \textsc{Gauss} zwar schon in seinem Manuskript nachträglich hinzugefügt. Da
das aber nur ganz vereinzelt und flüchtig geschehen ist und auch vielfach noch nicht die richtigen Werte
gegeben waren, wurden sie nicht mitabgedruckt. Auch sonst fanden sich einige Unrichtigkeiten,
die verbessert worden sind. \textsc{Gauss} hat hier offenbar seine Untersuchungen ziemlich plötzlich abgebrochen,
ohne seine Rechnungen bereits geprüft zu haben.
\begin{flushright}\textsc{Brendel}.\end{flushright}
\par\begin{center}\rule{.24\linewidth}{.4pt}\end{center}
\clearpage
\endgroup
